\documentclass[a4paper,11pt]{article}
\usepackage[utf8]{inputenc}
\usepackage[T1]{fontenc}
\usepackage[french]{babel}
\usepackage[left=2cm,right=2cm,top=2cm,bottom=2cm]{geometry}
\usepackage{xcolor}
\usepackage{hyperref}
\usepackage{fontawesome5}
\usepackage{parskip}

% --- Couleurs Thales ---
\definecolor{primary}{RGB}{0, 45, 114}
\hypersetup{colorlinks=true, urlcolor=primary}

\begin{document}

% --- EXPÉDITEUR ---
\noindent
\begin{minipage}[t]{0.5\textwidth}
    \textbf{Loïc Aron MBASSI EWOLO}\\
    Élève-Ingénieur Bac+5 -- Double diplôme ENSEA\\[3pt]
    \faMapMarker\ Cergy, France\\
    \faPhone\ (+33) 7 59 52 33 98\\
    \faEnvelope\ \href{mailto:loicaron.mbassiewolo@ensea.fr}{loicaron.mbassiewolo@ensea.fr}
\end{minipage}
\hfill
\begin{minipage}[t]{0.4\textwidth}
    \raggedleft
    Cergy, le 20 mars 2026
\end{minipage}

\vspace{1.2cm}

\hfill
\begin{minipage}[t]{0.5\textwidth}
    \raggedleft
    \textbf{Thales}\\
    Service Recrutement\\
    Élancourt, France
\end{minipage}

\vspace{1cm}

\noindent\textbf{Objet :} Candidature pour l'alternance Ingénieur IA pour l'Ingénierie Système Aéro (F/H)

\vspace{0.8cm}

Madame, Monsieur,

Générer automatiquement une architecture backend complète à partir de diagrammes UML grâce à un outil de parsing XML couplé à un LLM : c'est le projet open source que j'ai initié et qui compte aujourd'hui plus de 30 étoiles sur GitHub. En découvrant l'alternance Ingénieur IA pour l'Ingénierie Système Aéro à Élancourt, j'ai immédiatement vu le lien entre ma démarche (utiliser l'IA pour accélérer et fiabiliser l'ingénierie logicielle) et votre objectif (déployer des solutions IA pour accroître l'efficacité de vos architectes et ingénieurs systèmes aéro).

Concrètement, cet outil (Laravel API Generator) analyse les diagrammes UML exportés en XML depuis draw.io, calcule géométriquement les relations entre composants et génère modèles, contrôleurs, migrations et routes, avec une assistance LLM pour les cas complexes. Cette expérience m'a formé à la modélisation formelle, au parsing structurel et à l'intégration d'IA dans des processus d'ingénierie, trois compétences directement transposables à vos missions. En parallèle, le challenge CoHoMa (robotique autonome militaire, collaboration ENSEA/Armée Française) m'a plongé dans un contexte de défense : optimisation d'IA embarquée sur Nvidia Jetson, détection YOLOv10, SLAM, fusion Lidar 3D et caméra de profondeur, le tout avec des contraintes de fiabilité et de reproductibilité (Docker).

Mon projet académique de Fault Detection \& Fault-Tolerant Control sur drone (MATLAB/Simulink) complète ce profil sur le volet aéro : modélisation dynamique en conditions nominales et dégradées, détection/isolation de défauts par observateurs, et stratégies de contrôle tolérant aux pannes. Mon stage à l'INRIA Saclay (Institut Polytechnique de Paris, Avril-Août 2026) apporte la rigueur de la rédaction scientifique et du NLP à grande échelle avec BERT/Transformers, pertinente pour la documentation technique et l'analyse textuelle d'exigences que vous mentionnez.

Je suis disponible dès septembre 2026 pour intégrer votre site d'Élancourt. Contribuer à transformer les processus d'ingénierie système aéro chez Thales grâce à l'IA représente pour moi une opportunité concrète de conjuguer mes compétences en IA, en systèmes embarqués et en ingénierie logicielle au service d'enjeux de défense à fort impact. Je serais ravi d'en discuter lors d'un entretien.

Je vous prie d'agréer, Madame, Monsieur, l'expression de mes salutations distinguées.

\vspace{0.8cm}

\textbf{Loïc Aron MBASSI EWOLO}\\
\textit{Élève-Ingénieur ENSEA / Polytechnique Yaoundé}

\end{document}
